\section{Выводы по главе}

В теоретической главе был проведён комплексный анализ проблематики восстановления истории построения CAD-моделей с использованием методов машинного обучения. Основные выводы могут быть сформулированы следующим образом:

\begin{enumerate}
\item Актуальность задачи. Проблема «немой» геометрии, возникающая при импорте моделей из нейтральных форматов, является острой практической проблемой для промышленных предприятий. Существующие решения - как алгоритмические, так и коммерческие - не обеспечивают автоматизированного восстановления истории построения с приемлемым качеством.

\item Параметрическая природа CAD-моделей. Понимание структуры дерева построения, типов операций (выдавливание, вырезание, скругление, фаска) и их геометрических признаков является необходимым условием для разработки методов их распознавания.

\item Форматы представления данных. STEP обеспечивает точное B-Rep представление, необходимое для детерминированного вычисления параметров, но сложен для прямого использования в нейросетях. STL предоставляет унифицированное сеточное представление, удобное для ML, но теряет семантику и топологию. Оптимальным является комбинированное использование обоих форматов.

\item Обзор нейросетевых архитектур. Анализ существующих подходов к анализу 3D-геометрии показал, что MeshCNN является наиболее подходящей архитектурой для задачи сегментации CAD-моделей, поскольку она работает непосредственно с треугольными сетками, сохраняет топологическую информацию и использует инвариантные признаки.

\item Гибридный подход как перспективное направление. Как показывают работы Autodesk Research и результаты собственных экспериментов, прямое предсказание параметрических команд нейросетью сталкивается с фундаментальными ограничениями. Гибридный подход, в котором нейросеть отвечает за распознавание и локализацию (сегментацию), а детерминированные алгоритмы CAD-ядра - за точное вычисление параметров, является наиболее перспективным.

\item Открытые проблемы. Несмотря на прогресс, остаются нерешёнными вопросы: нехватка размеченных датасетов, проблема обобщения на сложную геометрию, чувствительность к численным параметрам, интерпретируемость и интеграция с существующими CAD-системами.
\end{enumerate}

Проведённый теоретический анализ создаёт необходимую основу для перехода к практической части работы, в которой будут рассмотрены конкретные архитектурные решения, реализация программного модуля и экспериментальные результаты.

